% Part I introduction (no \chapter heading; flows after the \part page).

\begin{draftblock}
\noindent In this part we are concerned with what economics has to say about \emph{agency itself}: what are reward, credit assignment and rational behaviour in the first place? This is not (yet) about aligning an AI to a principal, but about agency in general. The thesis of this part is that the market provides us not just a metaphor for reward, credit assignment and rational behaviour, but a theory of them.
\end{draftblock}

% Lifted (adapted) from "thesis or report structure.org":
When this research programme started, the motivating observation was that markets are \emph{naturally boundedly rational}: the efficient market hypothesis is best read as saying that markets are efficient \emph{conditional on the algorithmic information available}, which is exactly the property we want from a theory of bounded agents.

% Adapted from the year-2 report, P1.
\Cref{ch:market-rl} presents ``Market-based architectures in RL and beyond'': it introduces a new type of market algorithm in which the state itself is factored into several axes called ``goods'', allowing for greater specialization and parallelism than existing market-based RL algorithms; demonstrates that market algorithms can generalize neural networks, but with a more general learning algorithm that has potential to address challenges such as \emph{search}, \emph{dynamic scale} and \emph{complete feedback}; and proposes ways market algorithms may be applied in conjunction with Large Language Models.

\begin{draftblock}
\Cref{ch:agency} then develops the underlying theory: starting from the observation that \emph{reward is hyperstitional information} (a Bregman-divergence/mirror-ascent view of what a reward signal communicates to a learner), it argues that \emph{price recursion is the rational theory of reward} --- the recursive propagation of value through an economy of computations, of which backpropagation and TD-learning arise as special cases. The accompanying \texttt{marketgrad} library implements some ideas for using markets as an optimization algorithm on production graphs.
\end{draftblock}
